\documentclass[runningheads]{llncs}

% --------------------------------------------------
% Packages
% --------------------------------------------------
\usepackage{amsmath,amssymb,amsfonts}
\usepackage{graphicx}
\usepackage{booktabs}
\usepackage{array}
\usepackage{multirow}
\usepackage{tabularx}
\usepackage{siunitx}
\usepackage{url}
\usepackage{xurl}
\usepackage{float}
\usepackage{hyperref}

\urlstyle{same}

\title{AI-Driven Forensic Indexing:\\
A Behavioral Model for Chain Snatching Identification}
\titlerunning{AI-Driven Forensic Indexing}
\author{Anonymous}
\authorrunning{Anonymous}
\institute{Affiliation and email to be inserted}
\begin{document}

\maketitle

\begin{abstract}
Automated analysis of CCTV footage for chain-snatching incidents is
challenging due to the large volume of surveillance data, varying
visual conditions, and complex interactions between pedestrians and suspects. Existing approaches often address individual
tasks such as object detection, motion analysis, depth estimation, or
action recognition independently, limiting their effectiveness for
forensic investigation. While recent surveillance systems integrate multiple vision tasks, few jointly address forensic evidence preservation, computational efficiency, and domain-specific snatch signatures. This paper presents an AI-driven forensic
search framework that integrates these capabilities into a progressive
three-tier architecture. The framework combines low-level motion and
visual preprocessing, semantic object detection and spatial reasoning,
and higher-level behavioural analysis to identify chain-snatching-related
interactions and generate searchable forensic events. Particular
emphasis is placed on modelling suspect--victim interactions,
including pillion-rider behaviour, while preserving relevant evidence
throughout the processing pipeline. The framework was evaluated on
43 CCTV video sequences containing 35 annotated chain-snatching events
and 8 negative sequences. The complete system achieved
\textbf{100.00\% precision, 94.29\% recall, 97.06\% F1-score, and
95.35\% accuracy}. Progressive processing reduced the retained search
space to \textbf{2.38\%} of the original video frames, while
\textbf{33 of 35 ground-truth events} were preserved through the
complete pipeline. The forensic indexing component additionally
achieved \textbf{100\% search accuracy and completeness} with
sub-millisecond average query latency. The results demonstrate that
progressive visual and behavioural analysis can substantially reduce
CCTV investigation workload while retaining relevant forensic evidence
and enabling efficient retrieval of crime-related events.
\end{abstract}

\keywords{Chain Snatching Detection, Forensic Surveillance, Computer Vision,
Object Detection, Monocular Depth Estimation, Behavioral Recognition,
Forensic Indexing}

\section{Introduction}

The increasing deployment of CCTV systems has produced large volumes of
surveillance footage that are difficult to inspect manually. In
chain-snatching investigations, relevant events may occur within only a
few frames and involve complex interactions between victims, suspects,
and motorcycles. Frame-level object detection alone is therefore
insufficient for reliable forensic analysis.

Existing approaches often focus on individual surveillance tasks, such as
motion detection, object detection, depth estimation, tracking, or action
recognition~\cite{roy2018snatch}, and such approaches may pass large numbers of
irrelevant observations to downstream stages and provide limited
support for structured forensic retrieval, while some recent systems integrate
multiple vision capabilities~\cite{park2026foresea}. However, these
approaches are generally not designed specifically for chain-snatching
analysis with domain-specific snatch signatures. Furthermore, generic action
recognition does not explicitly model the spatial and temporal
relationships characteristic of chain-snatching incidents.

This paper presents an AI-driven forensic search framework that
progressively transforms CCTV footage into structured forensic events.
The framework integrates motion filtering, semantic object detection,
tracking, spatial relationship analysis, behavioural reasoning,
pose-based action recognition, evidence fusion, crime-specific
signature analysis, and forensic indexing. The progressive architecture
is designed to reduce the search space while preserving relevant
forensic evidence.

The framework is evaluated on a CCTV dataset containing annotated
chain-snatching and negative sequences. Experimental results demonstrate
the effectiveness of the proposed progressive architecture in reducing
the forensic search space while preserving relevant crime evidence and
supporting efficient CCTV-based chain-snatching investigation.

\section{Taxonomy of the Literature Survey}

To provide a structured overview of the reviewed literature, the surveyed techniques are categorized into three major research components: preprocessing, localization, and behavioral modeling. Each component is further divided into its respective subdomains based on the primary objective of the methods. Figure~\ref{fig:taxonomy} illustrates the taxonomy adopted in this survey.

\begin{figure}[!htbp]
    \centering
    \includegraphics[width=0.48\textwidth]{chain_snatching_taxonomy.png}
    \caption{Taxonomy of the literature surveyed for AI-driven chain-snatching detection.}
    \label{fig:taxonomy}
\end{figure}

\begin{table}[!t]
\centering
\scriptsize
\renewcommand{\arraystretch}{1.2}
\caption{Comparison of representative low-light image enhancement techniques for forensic surveillance.}
\label{tab:lowlight_comparison}

\begin{tabular}{|l|p{3.0cm}|p{3.5cm}|p{3.8cm}|p{2.5cm}|}
\hline
\textbf{Model} &
\textbf{Core Technique} &
\textbf{Strength} &
\textbf{Primary Limitation} &
\textbf{Suitability} \\ \hline

LIME~\cite{lime} &
Illumination Map Estimation &
Simple illumination correction &
High computational cost; prone to over-exposure under strong lighting &
Limited \\ \hline

Retinex-Net~\cite{retinexnet} &
Deep Retinex Decomposition &
Improves contrast and detail restoration &
Amplifies sensor noise under challenging lighting conditions &
Moderate \\ \hline

EnlightenGAN~\cite{enlightengan} &
GAN-based Image Enhancement &
Produces visually appealing enhanced images &
May generate hallucinated textures, compromising evidentiary reliability &
Low \\ \hline

\textbf{Zero-DCE}~\cite{zerodce} &
Zero-Reference Curve Estimation &
Non-generative, lightweight, and preserves structural details &
Performance may degrade under extremely severe illumination conditions &
High \\ \hline

CycleGAN~\cite{cyclegan} &
GAN-based Image Translation &
Effective style transfer without paired training &
May alter image content through hallucinated structures &
Low \\ \hline

\end{tabular}
\end{table}

\section{Subdomain Analysis \& Model Evolution}
\label{sec:subdomain}

\subsection{Low-Light Image Enhancement}

Continuous urban surveillance must operate under diverse illumination conditions, including nighttime environments where reduced visibility, sensor noise, motion blur, and low contrast can obscure critical visual cues such as hand movements and interactions between suspects, victims, and motorcycles. These conditions can degrade object detection and subsequent behavioral recognition. Furthermore, benchmark surveillance datasets such as UCF-Crime are predominantly composed of daytime footage, limiting the robustness of many existing computer vision models when deployed under low-light conditions. As a result, object detectors and subsequent behavioral recognition models experience substantial performance degradation due to the loss of discriminative image features. Therefore, low-light enhancement is an important preprocessing stage for improving visibility while preserving the structural and evidentiary information required for downstream forensic analysis.

\subsubsection{Comparative Analysis}

\paragraph{Illumination Map Estimation}
Classical techniques such as LIME estimate illumination maps to improve image luminance. However, they can introduce over-exposure when strong light sources such as street lamps or vehicle headlights are present~\cite{lime}.

\paragraph{Generative Models}
Unpaired image-to-image translation models such as EnlightenGAN improve perceptual quality but may introduce hallucinated textures or details. Such alterations are undesirable in forensic applications where preservation of the original visual evidence is essential~\cite{enlightengan}.

\paragraph{Non-Reference Curve Estimation}
Zero-Reference Deep Curve Estimation (Zero-DCE) formulates low-light enhancement as an image-specific curve estimation problem. A lightweight Deep Curve Estimation Network (DCE-Net) predicts pixel-wise enhancement curves that are iteratively applied to adjust image brightness while preserving local contrast and color information~\cite{zerodce}.

The light-enhancement curve is expressed as

\begin{equation}
L_{n}(x)=L_{n-1}(x)+\alpha_{n}L_{n-1}(x)
\left(1-L_{n-1}(x)\right),
\label{eq:zerodce}
\end{equation}

where $L_n(x)$ represents the enhanced pixel intensity after the $n^{\text{th}}$ iteration, $L_{n-1}(x)$ is the output of the previous iteration, and $\alpha_n$ controls the enhancement magnitude. Zero-DCE uses non-reference losses for spatial consistency, exposure control, color constancy, and illumination smoothness, eliminating the need for paired or unpaired training images~\cite{zerodce}. Its lightweight, non-generative formulation also avoids introducing synthesized details, making it suitable for forensic surveillance.


\subsection{Background Subtraction \& Temporal Filtering}

Background subtraction identifies foreground objects in video streams and reduces irrelevant background information before subsequent detection and behavioral analysis. This is particularly useful for chain-snatching detection, where relevant moving objects include pedestrians, suspects, victims, and vehicles. However, outdoor surveillance introduces challenges such as illumination changes, camera motion, shadows, moving backgrounds, adverse weather, and motion blur, which can reduce foreground localization accuracy.

\subsubsection{Comparative Analysis}

\paragraph{Classical Frame Differencing}
Frame differencing extracts foreground regions by subtracting consecutive frames. Although computationally efficient, it is sensitive to camera motion, illumination changes, and object speed, often producing fragmented masks and false detections~\cite{framediff}.

\paragraph{Gaussian Mixture Models}
Gaussian Mixture Models (GMM) model each pixel using a combination of Gaussian distributions and can adapt to gradual background changes. However, sudden illumination changes and complex outdoor scenes can reduce segmentation accuracy~\cite{gmm}.

\paragraph{Sample-Based Background Modeling}
Sample-based methods such as ViBe maintain previously observed pixel samples and provide efficient adaptation to changing backgrounds. However, stationary foreground objects may gradually become incorporated into the background and remain undetected~\cite{vibe}.

\paragraph{Adaptive Gaussian Mixture Models}
MOG2 extends GMM by adaptively estimating the number of Gaussian components for each pixel. This improves robustness to gradual illumination changes, shadows, and repetitive background motion while maintaining real-time performance. Its balance between segmentation quality and computational efficiency makes it suitable for continuous outdoor surveillance~\cite{zivkovic2004mog2}.


\subsection{Real-Time Object Detection}

Object detection localizes and classifies objects within individual video frames and provides the semantic information required for subsequent depth estimation and spatio-temporal analysis. In chain-snatching detection, accurate localization of pedestrians, motorcycles, and other road users is essential. However, motion blur, partial occlusion, varying object scales, crowded scenes, and adverse illumination require a balance between detection accuracy and computational efficiency~\cite{malhotra2020objectcomparison}.

\subsubsection{Comparative Analysis}

\paragraph{Two-Stage Object Detectors}
Two-stage detectors such as R-CNN and Faster R-CNN first generate candidate regions and subsequently classify and localize them. Although they generally provide accurate localization, their computational cost limits their suitability for real-time surveillance  where high-resolution videos need to be processed~\cite{fasterrcnn}.

\paragraph{Single-Stage Object Detectors}
Single-stage detectors such as SSD and the YOLO family perform detection in a single pass, providing faster inference and making them more suitable for continuous surveillance~\cite{ssd,yolo}. While SSD offers computational efficiency, detection performance can
vary for small or visually challenging objects. Whereas YOLO variants progressively improve the accuracy--speed trade-off and are widely used for real-time object detection~\cite{yolov8,yolov10,yolov11}.

\paragraph{YOLOv11}
YOLOv11 incorporates the C3k2 block and C2PSA attention module to improve feature extraction while maintaining computational efficiency~\cite{yolov11}. Its accuracy, parameter efficiency, and inference speed compared to its predecessors make it suitable for detecting pedestrians and two-wheeled vehicles in challenging surveillance footage using multi-scale feature extraction and attention
mechanisms.


\subsection{Monocular Depth Estimation}

A major challenge in CCTV-based crime analytics is 2D occlusion, where overlapping bounding boxes may incorrectly suggest physical interaction. For example, a pedestrian walking near a motorcycle may appear spatially close in image coordinates despite being physically separated. This can lead to erroneous interaction reasoning and false positives in chain-snatching detection.

\subsubsection{Comparative Analysis}

\paragraph{Stereo Vision and Active Depth Sensors}
Traditional depth estimation methods use stereo configurations or active sensing approaches such as LiDAR and Time-of-Flight (ToF). Although these methods provide accurate depth observations, they require additional hardware, calibration, and deployment cost, limiting their suitability for large-scale urban surveillance.

\paragraph{Classical Monocular Depth Estimation}
Deep monocular approaches eliminate the need for specialized depth sensors. MiDaS demonstrated strong cross-dataset relative-depth estimation, but its computational requirements and reduced boundary precision in complex scenes can limit its suitability for real-time surveillance~\cite{monovit}.

\paragraph{Foundation-Based Depth Estimation}
Depth-Anything V2~\cite{yang2024depth} uses a Vision Transformer foundation model trained on large-scale synthetic and real-world data. It provides robust relative-depth estimation, sharper object boundaries, and improved generalization across environments, making it suitable for estimating spatial relationships between objects in complex urban surveillance scenes.


\subsection{Pose Estimation and Tier Integration}

Following object detection and monocular depth estimation, human pose estimation provides an intermediate representation of detected persons through skeletal keypoints. Compared with appearance-based features, skeletal representations are less sensitive to background clutter and illumination while retaining information about body posture and motion. Pose estimation therefore provides a bridge between lower-level visual analysis and higher-level behavioral recognition.

\subsubsection{MediaPipe Pose as the Bridge Module}

MediaPipe Pose is used as the intermediate module connecting object detection and monocular depth estimation with behavioral recognition~\cite{mediapipe}. Its heatmap-based keypoint estimation provides a skeletal representation that can retain useful information under partial occlusion, which is important during rapid chain-snatching interactions.

MediaPipe Pose also incorporates temporal tracking, allowing consecutive frames to be processed efficiently without repeatedly performing full-scale pose estimation. This reduces computational cost while maintaining short-term consistency in skeletal motion. In addition to 2D joint locations, its depth estimates under a weak-perspective camera model can be combined with the global depth estimates from Depth-Anything V2~\cite{yang2024depth} to form a more informative 3D skeleton for the ST-GCN model~\cite{yan2018stgcn}.


\subsection{Spatio-Temporal Behavioral Recognition}

\begin{table}[!t]
\centering
\scriptsize
\renewcommand{\arraystretch}{1.2}
\caption{Quantitative comparison of representative behavior recognition models for forensic surveillance.}
\label{tab:behavior_quantitative}

\begin{tabular}{|l|c|c|c|}
\hline
\textbf{Metric} &
\textbf{ViViT~\cite{liu2022video}} &
\textbf{2s-AGCN~\cite{shi2019twostreamagcn}} &
\textbf{ST-GCN~\cite{yan2018stgcn}} \\ \hline

Top-1 Accuracy (NTU-120) (\%) $\uparrow$ &
89.1 &
\textbf{91.2} &
89.9 \\ \hline

GFLOPs $\downarrow$ &
50.0 &
18.4 &
\textbf{16.2} \\ \hline

Parameters (Millions) $\downarrow$ &
100+ &
6.9 &
\textbf{3.1} \\ \hline

Skeleton-Based Representation &
No &
Yes &
Yes \\ \hline

Real-Time Deployment &
Moderate &
High &
\textbf{High} \\ \hline

Interpretability &
Low &
High &
\textbf{High} \\ \hline

\end{tabular}
\end{table}

\subsubsection{Overview}

The final stage distinguishes legitimate pedestrian--vehicle interactions from visually similar chain-snatching events. Object detection and depth estimation can identify entities and estimate their spatial relationships, but they cannot determine whether an interaction represents criminal behavior. Behavioral recognition therefore models body configuration and motion across successive frames to identify fine-grained spatio-temporal patterns associated with chain snatching.

\subsubsection{Comparative Analysis}

\paragraph{CNN-LSTM}
Early approaches combined CNNs for spatial feature extraction with LSTMs for temporal modelling. Although effective for simple activities, these methods can be sensitive to background appearance, illumination, and viewpoint changes~\cite{cnnlstm}.

\paragraph{Two-Stream Networks}
Two-stream architectures process RGB information and optical flow separately, explicitly modelling motion. However, computing dense optical flow for every frame introduces substantial computational cost for continuous surveillance~\cite{twostreamlstm}.

\paragraph{Vision Transformer-Based Models}
Vision Transformer models such as ViViT use self-attention to model long-range temporal dependencies but require comparatively large computational and memory resources, limiting their deployment in resource-constrained surveillance systems~\cite{liu2022video}.

\paragraph{Graph-Based Skeleton Recognition}
Graph-based methods represent the human body as a graph whose nodes correspond to joints and edges represent anatomical connections. ST-GCN learns spatial configurations and temporal motion patterns directly from skeletal data, providing greater robustness to illumination, background, clothing, and viewpoint variations. Adaptive approaches such as 2s-AGCN can improve recognition accuracy by learning the graph structure, but introduce additional computational cost~\cite{shi2019twostreamagcn,yan2018stgcn}.

ST-GCN therefore provides a suitable balance between computational cost, accuracy, and real-time deployment for the proposed system. Combined with MediaPipe Pose and Depth-Anything V2, it enables depth-aware spatio-temporal analysis of pedestrian--motorcycle interactions and supports discrimination between chain-snatching and visually similar legitimate interactions~\cite{yan2018stgcn}.

\section{Research Gaps}
\label{sec:gaps}

The comparative analysis in Section~3 identifies several limitations that
hinder the application of existing surveillance approaches to automated
chain-snatching detection. Although substantial progress has been made in
low-light enhancement, object detection, depth estimation, and behavioural
recognition, existing solutions do not adequately address the requirements
of complex CCTV environments. The main research gaps identified from the
literature are described below.

\subsection{Pillion-Rider Behavioral Gap}

Existing action recognition approaches primarily target generic human
behaviours such as walking, running, fighting, and falling. Although
graph-based approaches such as ST-GCN effectively model human motion
patterns~\cite{yan2018stgcn}, fine-grained interactions involving pedestrians,
motorcycles, and rapid ``grab-and-accelerate'' motion in chain-snatching
scenarios remain insufficiently modelled~\cite{roy2018snatch}. Consequently,
existing action datasets and recognition models provide limited support for
identifying pillion-rider chain-snatching behaviour.

\subsection{Nighttime Reliability Gap}

Existing surveillance datasets are predominantly biased toward daytime
conditions, while nighttime footage suffers from poor illumination, motion
blur, and reduced sensor sensitivity. Traditional enhancement methods may
struggle under these conditions, whereas generative approaches can introduce
artifacts that are not present in the original footage
\cite{lime,enlightengan,zerodce,retinexnet}. This creates a need for
efficient non-generative enhancement methods that improve visual information
while preserving the original evidence.

\subsection{Edge Resource Bottleneck}

Increasing recognition capability is often accompanied by higher
computational and memory requirements. Computationally intensive
video-recognition architectures can therefore limit deployment on
resource-constrained edge platforms~\cite{liu2022video,feichtenhofer2019slowfast}.
This motivates progressive and resource-efficient architectures for forensic
video analysis that reduce unnecessary computation without relying entirely
on cloud-based processing.

\subsection{2D Perspective Ambiguity}

Reasoning about interactions using only two-dimensional image coordinates can
produce incorrect proximity relationships because objects at different
physical depths may overlap or appear close in the image plane. This
perspective ambiguity can lead to false interaction events
\cite{monovit,yang2024depth}. Incorporating monocular depth
estimation can therefore provide additional spatial information for
validating interactions in surveillance footage.

\section{Proposed Methodology}
\label{sec:methodology}

\subsection{System Overview}

The proposed approach is three-stage hierarchical in nature where deep learning models of increasing complexity are applied sequentially only when needed. Firstly, motion filtering is applied to discard stationary frames, followed by spatial localization and finally performing spatio-temporal behavioural analysis. This approach reduces the overall computational cost while retaining the forensic accuracy.

\subsection{Motion Triage and Illumination Normalization}
\subsubsection{Background Subtraction using MOG2}

Motion triage is performed using the MOG2 (Mixture of Gaussians) background subtractor \cite{zivkovic2004mog2}, which models each pixel as an adaptive mixture of Gaussians to distinguish foreground objects from the background under changing scene conditions. Its adaptive learning rate accommodates gradual illumination changes caused by weather conditions, while OpenCV's shadow detection masks shadowed regions before further processing.

In the proposed framework, MOG2 is configured with a history length of 500 frames, a variance threshold of 16.0, and shadow detection enabled. For each incoming frame, a foreground mask is generated and the number of non-zero foreground pixels is computed. Frames are retained only when more than 5,000 foreground pixels are detected; otherwise, they are discarded without invoking subsequent processing stages. Frames satisfying this criterion are forwarded to the object detection stage, thereby reducing unnecessary computation.

\subsubsection{Low-Light Enhancement using Zero-DCE}

Frames containing foreground activity are conditionally passed to the
Zero-DCE enhancement network \cite{zerodce} when poor illumination is
detected. The illumination level is estimated from the mean brightness of
the motion-triaged region in the HSV colour space,

\begin{equation}
\bar{L}(f)=\operatorname{mean}\left(V_{\mathrm{HSV}}(f)\right),
\end{equation}

where $V_{\mathrm{HSV}}(f)$ denotes the value (brightness) channel of frame
$f$. Enhancement is applied only when

\begin{equation}
\bar{L}(f)<\tau_{lux},
\end{equation}

where $\tau_{lux}$ is a tunable illumination threshold selected according
to the deployment environment. Thus, well-illuminated frames are retained
without enhancement, reducing unnecessary computation.

Zero-DCE performs enhancement by estimating pixel-wise adjustment curves
rather than directly generating a new image:

\begin{equation}
L_n(x)=L_{n-1}(x)+\alpha_nL_{n-1}(x)
\left(1-L_{n-1}(x)\right),
\end{equation}

where $L_n(x)$ represents the enhanced pixel intensity after the
$n^{\text{th}}$ iteration and $\alpha_n$ controls the enhancement strength.
The process progressively brightens dark regions while limiting
over-exposure in already illuminated regions.

The Zero-DCE network uses the lightweight fully convolutional DCE-Net,
consisting of seven $3\times3$ convolutional layers with symmetric skip
connections~\cite{zerodce}. Unlike encoder-decoder architectures, DCE-Net
operates at the input resolution without spatial pooling or upsampling,
enabling real-time inference exceeding 500 FPS. Its final convolutional
layer predicts eight sets of RGB enhancement-curve parameters, constrained
using a \textit{tanh} activation and applied iteratively over eight
enhancement stages.

\subsection{Object Detection, Depth and Tracking}

\subsubsection{Object Detection using YOLO11n}

Frames satisfying the motion-triage criterion are processed using
YOLO11n \cite{yolov11} to localize objects relevant to subsequent interaction analysis.
The detector produces bounding boxes together with the corresponding
class labels and confidence scores. GPU acceleration is used when
available, with CPU execution provided as a fallback.

For each detected object, the bounding-box geometry is represented using
its corner coordinates $(x_1,y_1,x_2,y_2)$, from which the box width,
height, area, and centroid are derived. The centroid is computed as

\begin{equation}
c_x = \frac{x_1+x_2}{2}, \qquad
c_y = \frac{y_1+y_2}{2}.
\end{equation}

The bounding-box dimensions and area are given by

\begin{equation}
w=x_2-x_1,\qquad
h=y_2-y_1,\qquad
A=w h.
\end{equation}

These spatial attributes provide the geometric representation required
by the subsequent semantic filtering, depth estimation, tracking, and
interaction-analysis stages.

\subsubsection{Semantic Filtering}

A semantic filtering layer is applied to the detections from the object detector in an effort to limit downstream processing to objects relevant to chain-snatching analysis. The filter only keeps the object categories that are relevant for crime: person, motorcycle, bicycle, car, bus and truck. This constraint cuts down irrelevant background objects and environmental clutter irrelevant to the subsequent interaction analysis.

Apart from the semantic class constraint, the detections are required to satisfy a confidence threshold above 0.35 and a minimum bounding-box area criterion. Detections that do not satisfy these constraints are discarded before they are fed to the next depth estimation and tracking stages. The filtered detections that result become the set of candidate entities for downstream spatial, temporal and behavioral analysis.

\subsubsection{Monocular Depth Estimation using Depth-Anything V2}

After object detection, the proposed framework uses Depth-Anything V2
\cite{yang2024depth} to estimate a dense relative depth map for each
illumination-normalized surveillance frame. The model uses a DINOv2 Vision
Transformer (ViT) backbone for semantic feature extraction and a Dense
Prediction Transformer (DPT) decoder to generate high-resolution depth maps
from a single RGB image.

The estimated depth map complements the 2D object detections produced by
YOLOv11 by providing additional 3D scene information. This helps distinguish
physical interactions from apparent intersections caused by perspective
projection.

Depth estimation operates in parallel with object detection so that both
spatial localization and relative depth information are available for
interaction validation. Candidate suspect--victim pairs satisfying the
spatial proximity and depth-consistency criteria are then forwarded. 

\subsubsection{Multi-Object Tracking}

Temporal correspondence between detections across consecutive
frames is established by a multi-object tracking layer. The tracker
associates bounding boxes over time and assigns persistent track identities
to individual entities, which allows their trajectories to be maintained across
successive observations and short-term occlusions.

The tracking mechanism combines Kalman-filter-based state prediction with
intersection-over-union (IoU) based data association, following the matching
principles used in modern tracking approaches such as ByteTrack~\cite{bytetrack} and
BoT-SORT. Each active track maintains a persistent identity together with
its historical centroid and bounding-box observations. These trajectories
provide the temporal information required by the subsequent motion-feature
and relationship-analysis stages.

\subsubsection{Motion Feature Extraction}

The temporal trajectories generated by the tracking stage are further
processed to derive motion descriptors for each tracked entity. The
centroid positions of a track across consecutive frames are used to
estimate its instantaneous velocity, average speed over a sliding temporal
window, acceleration, total displacement, and movement direction.

The resulting kinematic descriptors are associated with the corresponding
tracks and provide quantitative information about the movement of each
entity. These features are subsequently used by the relationship and
interaction-analysis stages to characterize relative motion between
tracked entities.

\subsection{Interaction and Behaviour Analysis}

\subsubsection{Relationship Analysis}

To characterize interactions between tracked entities, the framework
performs pairwise spatial and motion analysis between relevant person and
vehicle tracks. For each candidate pair, the Euclidean distance between
their image-space centroids is calculated as

\begin{equation}
d(p_1,p_2)=
\sqrt{(x_2-x_1)^2+(y_2-y_1)^2}.
\end{equation}

The framework additionally evaluates the relative closing speed and
spatial orientation between the entities. A pair is considered for
relationship analysis when its centroid distance falls within a
150-pixel threshold. These measurements provide a representation of the
spatial and relative-motion relationship between tracked entities and
are subsequently used by the interaction-management stage to identify
persistent interactions.

\subsubsection{Interaction Management}

Relationship observations generated in the preceding stage are maintained
over time using a finite-state interaction manager. The manager tracks
person--vehicle pairs across consecutive frames and maintains their
interaction lifecycle, allowing temporary separations or occlusions to be
distinguished from interaction termination.

An interaction is initialized in the \textit{NEW} state when a person--vehicle
pair satisfies the spatial proximity criterion and transitions to
\textit{ACTIVE} after a second consecutive observation. If the pair
temporarily disappears or separates, it enters the \textit{LINGERING} state,
allowing the relationship to persist through short-term observation gaps.
After 30 inactive frames, it transitions to \textit{ENDED} and subsequently
to \textit{ARCHIVED}. A lingering interaction returns to \textit{ACTIVE} when
the corresponding pair reappears.

For each active interaction, temporal, spatial, and motion histories are
maintained, including current, minimum, maximum, and running-average
distances, centroid positions, speeds, and movement headings. Relative
kinematic quantities include closing velocity, acceleration, heading
difference, and trajectory cosine similarity:

\begin{equation}
S_{\mathrm{traj}} =
\frac{\mathbf{d}_{p}\cdot\mathbf{d}_{v}}
{\|\mathbf{d}_{p}\|\|\mathbf{d}_{v}\|},
\end{equation}

where $\mathbf{d}_{p}$ and $\mathbf{d}_{v}$ denote the displacement vectors
of the person and vehicle, respectively.

The manager also computes a normalized interaction confidence score
$C\in[0,1]$ from proximity, interaction duration, and relative closing
speed:

\begin{equation}
C =
\min\left(1.0,\,
0.5S_{\mathrm{proximity}}
+0.3S_{\mathrm{duration}}
+0.2S_{\mathrm{velocity}}\right).
\end{equation}

The proximity component increases as distance decreases, while the duration
component saturates at 60 frames and the velocity component increases with
closing speed. These interaction histories and confidence estimates are
retained as structured evidence for subsequent stages.

\subsubsection{Behaviour Intelligence Layer}

The interaction histories generated by the preceding stages are converted
into discrete behavioural primitives using a rule-based behaviour
intelligence layer. At each frame, active interactions are evaluated using
spatial proximity, relative velocity, acceleration, heading, trajectory
similarity, and motion changes. Each satisfied rule generates a
time-stamped behavioural primitive, representing the evolution of an
interaction as a chronological sequence rather than isolated frame-level
observations.

The behavioural rules capture phases and motion patterns including
approaching, moving away, following, parallel motion, close interaction,
stationary interaction, rapid acceleration or deceleration, sudden
direction changes, and rapid separation. Approaching is identified from a
negative relative velocity below $-2.0$ pixels per frame, while close
interaction is associated with an inter-entity distance below 80 pixels.
Trajectory following is characterized by a cosine similarity above 0.70,
rapid acceleration by acceleration above 3.0 pixels per frame$^2$, and rapid
separation by relative velocity above 5.0 pixels per frame.

The resulting primitives are organized chronologically into a behaviour
timeline, providing a structured representation of the interaction sequence.
This representation is subsequently used by the behaviour graph and
multimodal fusion stages to derive higher-level behavioural
evidence~\cite{yan2018stgcn}.

\subsubsection{Behaviour Graph Modeling}

The behavioural primitives generated by the behaviour intelligence layer are
organized into a spatio-temporal directed graph that represents the evolution
of an interaction. The graph is constructed incrementally on a frame-by-frame
basis by evaluating the active interaction and selecting the highest-priority
behavioural pattern among the detected candidates.

Each behavioural pattern is represented by a graph node containing its
pattern type, temporal extent, confidence, priority, supporting primitives,
and associated motion and spatial evidence. When the primary pattern changes,
a new node is created; otherwise, the existing node is updated with the latest
end frame, duration, confidence, and kinematic statistics. This preserves both
persistent behavioural states and transitions between interaction phases.

Directed temporal edges are created when the active pattern transitions to a
different pattern. Each edge records the source and target patterns, transition
frame, video timestamp, and observed transition condition, thereby preserving
the temporal ordering of behavioural patterns.

Pattern detection considers spatial, kinematic, alignment, and temporal
criteria, including proximity, approach and separation velocity,
acceleration, stationary speed, heading alignment, trajectory similarity,
and minimum pattern duration. These criteria combine the spatial and motion
evidence accumulated by the preceding interaction stages to construct a
structured representation of behavioural evolution.

The resulting behaviour graph constitutes \textit{Stream A} of the
proposed framework. It provides high-level spatio-temporal behavioural
evidence that is subsequently combined with the pose-based action
recognition evidence generated by Stream B.

\subsubsection{ROI Selection and Preparation}

To provide stable spatial inputs for pose estimation, the framework extracts
regions of interest (ROIs) corresponding to the person involved in an active
interaction. The ROI is obtained from the tracked person's frame-wise
bounding boxes, thereby restricting subsequent pose analysis to the relevant
participant rather than unrelated subjects in the scene.

Missing bounding-box observations caused by temporary occlusions are handled
using linear interpolation for gaps of up to 15 frames. Unphysical bounding
box variations are identified using the median bounding-box area and its
standard deviation. Observations whose area deviates from the median by more
than three standard deviations are treated as outliers and replaced using
the previous valid bounding box.

Temporal jitter is subsequently reduced using Exponential Moving Average
(EMA) smoothing. For a bounding-box vector
$B_t=[x_1,y_1,x_2,y_2]^T$, the smoothed coordinates are computed as

\begin{equation}
B_t^{\mathrm{smooth}}
=
\alpha B_t^{\mathrm{raw}}
+
(1-\alpha)B_{t-1}^{\mathrm{smooth}},
\end{equation}

where $\alpha=0.4$.

To retain relevant contextual information, including extended body parts
and nearby vehicle-contact regions, the smoothed bounding boxes are
expanded by 25\% in width and height. The resulting ROI sequences are then
evaluated using sequence completeness, missing-frame ratio, bounding-box
stability, and temporal frame coverage. ROIs are accepted only when they
satisfy the corresponding quality criteria, thereby removing unreliable
crops before pose estimation.

The accepted ROIs are extracted from the corresponding video frames and
organized into interaction-specific sequences, typically covering
30--60 frames. The resulting interpolated, outlier-filtered, smoothed,
expanded, and quality-validated ROI sequences are forwarded to the
pose-estimation stage.

\subsection{Pose-Based Action Recognition}
\subsubsection{Pose Estimation}

The interaction-specific ROI sequences are processed by a human pose
estimation layer to extract frame-wise skeletal keypoints. The baseline
implementation uses MediaPipe Pose (BlazePose) to estimate anatomical
landmarks from each person ROI. Since MediaPipe Pose provides 33 body
landmarks, the detected landmarks are mapped to a standardized
17-joint COCO skeleton topology to provide a consistent representation
for the subsequent skeleton-based action recognition stage.

To improve temporal consistency, missing joint observations caused by
temporary occlusions are handled using linear interpolation for gaps of
up to 10 frames. Low-confidence joint estimates below a minimum joint
confidence threshold of 0.30 are discarded. 

A pose quality score is additionally computed from the mean joint
confidence and the proportion of valid joints, providing a measure of
the reliability of each estimated pose sequence. Low-quality pose
observations can therefore be excluded before skeleton sequence
construction.

\subsubsection{Skeleton Sequence Construction}

The temporally smoothed pose keypoints are transformed into normalized
skeleton sequences for graph-based action recognition. To reduce the
effect of variations in image position and subject location, the
baseline representation uses hip-centred normalisation, expressing
joint coordinates relative to the subject's hip region.

The normalized keypoints are organized into fixed-length temporal
sequences of $T=30$ frames. The resulting representation follows the
NCTVM tensor convention, where $N$ denotes the batch dimension, $C$ the
coordinate channels, $T$ the temporal dimension, $V$ the number of
skeletal joints, and $M$ the number of persons. For the proposed
implementation, the resulting tensor has the form

\begin{equation}
(1,C,T,V,1)=(1,2,30,17,1),
\end{equation}

where the two coordinate channels represent the 2D joint coordinates,
17 joints are used to represent the human skeleton, and a single
subject is represented in each sequence.

The resulting normalized skeleton tensors provide a consistent
spatio-temporal representation for the subsequent ST-GCN action
recognition stage.

\subsubsection{ST-GCN Action Recognition}

The normalized skeleton sequences are processed using a
Spatial-Temporal Graph Convolutional Network (ST-GCN)~\cite{stgcn} to
recognize human actions from their joint-level motion patterns. The
human body is represented as a graph in which the 17 skeletal joints
form the nodes and anatomical connections define the spatial edges.
Temporal connections between the same joints across consecutive frames
extend this representation into a spatio-temporal graph.

The spatial graph convolution operates on the normalized adjacency
matrix of the COCO-17 skeleton topology, while temporal convolution
models the evolution of joint configurations across the $T=30$ frame
sequence. Each graph-convolution block is followed by batch
normalization and a ReLU activation. Global average pooling is then
applied before a fully connected classification layer produces
probabilities over the ten action classes.

The action taxonomy consists of \textit{Walking}, \textit{Standing},
\textit{Running}, \textit{Approaching}, \textit{Reaching},
\textit{Grabbing}, \textit{Pulling}, \textit{Turning}, \textit{Falling},
and \textit{Unknown}. These classes provide domain-specific motion
representations for distinguishing ordinary pedestrian activity from
motion patterns associated with chain-snatching interactions.

The resulting action predictions constitute \textit{Stream B} of the
behavioural analysis framework. Unlike the rule-based spatio-temporal
behaviour graph in Stream A, this branch provides learned skeleton-based
action evidence. The two streams are subsequently combined by the
behaviour fusion stage to obtain a unified behavioural assessment.

Each inference produces the predicted action, its confidence, class
probabilities, top-$k$ predictions, and inference latency, together with
the associated sequence, interaction, and track identifiers. This
structured output allows the action evidence to be incorporated into the
subsequent behavioural fusion and forensic event-generation stages.

\subsection{Multimodal Behaviour Fusion}

\subsubsection{Behaviour Fusion}

The behavioural evidence generated by the two complementary analysis
streams is combined to obtain a unified interaction assessment. Stream A
provides explicit spatio-temporal evidence through the behaviour graph,
while Stream B provides learned human-action evidence through the
ST-GCN. Their combination reduces reliance on either rule-based
interaction reasoning or skeleton-based action recognition alone.

The primary fusion mechanism uses a weighted confidence formulation:

\begin{equation}
C_{\mathrm{fused}} =
w_A C_{\mathrm{graph}} +
w_B C_{\mathrm{action}},
\end{equation}

where $C_{\mathrm{graph}}$ and $C_{\mathrm{action}}$ denote the
confidence values obtained from the behaviour-graph and action-
recognition streams, respectively, and $w_A$ and $w_B$ represent their
corresponding fusion weights.

The resulting fused interaction representation combines the temporal
interaction structure, spatial and motion evidence, and human action
predictions into a common behavioural assessment. This fused evidence is
passed to the subsequent snatch-signature stage for crime-specific
reasoning.

\subsection{Crime-Specific Snatch Signature Reasoning}

\subsubsection{Snatch Signature Matching}

The fused behavioural evidence is evaluated by a crime-specific signature
engine to determine whether an interaction is consistent with a
chain-snatching event. Rather than directly classifying an entire video
sequence as a crime, the engine evaluates the observed interaction
against predefined behavioural signatures corresponding to different
snatching scenarios.

The framework defines separate signatures for motorcycle-based and
pedestrian-based snatching scenarios. Each signature evaluates the
available behavioural evidence using weighted criteria derived from the
interaction dynamics and behavioural patterns. The resulting signature
score is normalized to the interval $[0,1]$ and represents the degree of
agreement between the observed interaction and the corresponding
snatching pattern.

A score of $S \geq 0.70$ is treated as a high-confidence signature
match. In addition to the numerical score, the signature engine
generates an evidence-based explanation indicating which signature
conditions were satisfied or not satisfied. This provides an
interpretable decision layer between multimodal behavioural evidence and
the final forensic event record.

\subsubsection{Forensic Event Indexing and Retrieval}

The final stage converts the crime-specific assessments generated by the
snatch signature engine into structured forensic event records, builds sub-millisecond (latency remains below 0.1~ms) $O(1)$   inverted search indices. Each
event retains the decision and the associated behavioural and interaction
evidence, enabling detected incidents to be searched and retrieved
without repeatedly processing the original video data.

A multi-attribute inverted indexing scheme is used to organize forensic
events according to attributes including event identifier, video
identifier, decision, associated track identities, behavioural patterns,
detected actions, and descriptive tags. Separate indices over these
attributes enable direct retrieval of events matching specific forensic
queries.

In addition to the searchable index, the framework preserves associated
evidence artefacts, including structured event records, thumbnails, and
annotated video clips. The resulting outputs provide investigators with
both searchable metadata and the corresponding visual evidence required
for subsequent forensic examination.

\section{Experimental Results and Discussion}
\subsection{Experimental Setup}

\subsubsection{Dataset}

The proposed framework was evaluated on 43 real-world CCTV video
sequences collected from the Snatch 1.0 benchmark and supplementary
surveillance footage. The dataset contains 96,926 frames corresponding
to 3,877.05~s (64.62~min) of video. Of the 43 sequences, 35 contain
chain-snatching events and 8 represent normal surveillance footage used
as negative controls. All 43 videos were successfully decoded, with no
corrupted or unreadable files identified. The dataset contains
heterogeneous surveillance conditions, including multiple spatial
resolutions and frame rates, predominantly static outdoor cameras, and
daytime, mixed-lighting, and night-time recordings.

\subsubsection{Ground Truth and Evaluation Protocol}

Binary video-level ground-truth labels were available for all 43
sequences, comprising 35 positive chain-snatching instances and 8
negative surveillance sequences. For positive sequences, temporal
interaction boundaries were defined for the relevant snatching events.
The evaluation additionally incorporated annotated action categories
including Walking, Standing, Approaching, Reaching, Grabbing, Pulling,
and Escaping, together with interaction categories representing normal
and crime-related encounters.

The evaluation was designed to assess both forensic detection
effectiveness and computational efficiency. Event-level detection
performance was measured using precision, recall, F1-score, accuracy,
false-positive rate, and false-negative rate. Pipeline efficiency was
evaluated using processing latency, throughput, and progressive
search-space reduction. Stage-wise evidence preservation was measured
to determine whether early filtering stages discarded ground-truth
crime events. Architectural contribution was further examined through
configuration comparisons and component-level ablation experiments.

\subsubsection{Computational Environment}

The experiments were executed on a CPU--GPU computing environment using
an 8-core processor and NVIDIA CUDA acceleration. During pipeline
execution, average CPU utilization was 22.4\%, resident memory usage
was 824.5~MB, GPU VRAM allocation was approximately 850~MB, and average
GPU utilization was 18.5\%.

\subsection{Motion Detection Performance Evaluation}
This section provides the results of the experiments performed to evaluate the motion detection stage of the proposed AI system for chain-snatching detection. The goal of the experiments was to benchmark the keyframe selection stage of various classical motion detection algorithms. The evaluation criteria included the algorithm’s computational efficiency, its ability to filter irrelevant information, and its capacity to preserve motion continuity to enable effective behavioural analysis. The experiments were performed on the Snatch 1.0 CCTV dataset, which comprises 43 surveillance videos, yielding 215 benchmark runs in total.

\subsection{Motion Detection Performance Evaluation}

The motion detection stage was evaluated to benchmark the keyframe selection
capability of five approaches: Baseline (No Filtering), Frame Difference,
Gaussian Mixture Model (GMM), K-Nearest Neighbour (KNN), and Mixture of
Gaussians Version 2 (MOG2). Each detector processed the 43 videos in the
Snatch 1.0 CCTV dataset using identical parameter settings, resulting in
215 benchmark runs. The evaluation considered three criteria:
\textbf{processing speed}, \textbf{frame reduction}, and
\textbf{motion continuity}.

\subsubsection{Motion Detection Benchmarking}

For each video, a binary foreground mask was generated for every frame and
used to identify motion regions. Processing speed was measured in frames per
second (FPS), frame reduction quantified the percentage of frames discarded,
and motion continuity measured the preservation of uninterrupted motion
sequences required for subsequent behavioural analysis.

\begin{figure}[htbp]
    \centering
    \includegraphics[width=0.5\textwidth]{fps_comparison.png}
    \caption{Average Processing Speed (FPS) Comparison of Motion Detection Algorithms}
    \label{fig:fps}
\end{figure}

Figure~\ref{fig:fps} shows the processing-speed comparison. Frame Difference
achieves the highest throughput among the filtering methods at 444.82 FPS,
followed by MOG2 (144.18 FPS), GMM (114.16 FPS), and KNN (90.77 FPS).
However, computational speed alone does not indicate the suitability of a
motion detector for behavioural analysis.

\begin{table}[htbp]
\centering
\scriptsize
\setlength{\tabcolsep}{4pt}
\renewcommand{\arraystretch}{0.9}
\caption{Overall Motion Detection Benchmark Results}
\label{tab:benchmark}

\begin{tabular}{lccc}
\hline
\textbf{Detector} & \textbf{Reduction (\%)} & \textbf{FPS} & \textbf{Continuity} \\
\hline
Baseline          & 0.00  & 1001.35 & 0.995 \\
Frame Difference  & 50.83 & 444.82  & 0.610 \\
GMM               & 34.73 & 114.16  & 0.851 \\
KNN               & 4.00  & 90.77   & 0.992 \\
MOG2              & 5.17  & 144.18  & 0.985 \\
\hline
\end{tabular}
\end{table}

Table~\ref{tab:benchmark} shows that Frame Difference provides the greatest
frame reduction (50.83\%) and highest throughput among the filtering methods,
but its low continuity score (0.610) indicates substantial fragmentation of
motion sequences. GMM provides greater reduction than MOG2 but with lower
speed and continuity, while KNN preserves continuity but achieves only
4.00\% reduction at 90.77 FPS. MOG2 provides a balanced combination of
\textbf{frame reduction}, \textbf{processing speed}, and
\textbf{motion continuity}, achieving 144.18 FPS and a continuity score of
0.985. It was therefore selected for the proposed preprocessing stage.

\subsubsection{Motion Area Threshold Analysis}

The effect of the motion-area threshold was evaluated using threshold values
ranging from 0\% to 10\%. Frames were retained only when detected foreground
motion exceeded the selected threshold, and the resulting frame reduction
and motion continuity were averaged across the dataset.

\begin{figure}[!t]
\centering
\includegraphics[width=0.92\textwidth]{motion_area_threshold.png}
\caption{Motion Area Threshold Sensitivity Analysis}
\label{fig:threshold}
\end{figure}

Figure~\ref{fig:threshold} shows that increasing the threshold generally
increases frame reduction by discarding more low-motion frames, but can also
decrease motion continuity by removing subtle movements. MOG2 maintains
relatively stable continuity across moderate threshold values while
providing useful frame reduction, supporting its selection for the proposed
framework.

\subsection{Pipeline Bottleneck and Latency Analysis}

The computational cost of the proposed pipeline was analyzed by measuring
the total runtime and per-frame latency of the major processing stages. The results are summarized in Table~\ref{tab:pipeline_latency}. 
The results show that the computational burden is concentrated in the
object detection and tracking stages.

The tracking stage represents the largest contribution to the measured
runtime, accounting for 47.36\% of the total execution time, with an
average latency of 68.21~ms per frame. YOLO object detection is the
second-largest contributor, accounting for 41.57\% of the runtime and
requiring an average of 54.82~ms per frame. Together, these two stages
account for 88.93\% of the measured pipeline runtime, identifying them
as the principal computational bottlenecks. This runtime distribution is illustrated in
Fig.~\ref{fig:runtime}.

Motion filtering contributes 10.90\% of the total runtime, with an
average latency of 13.99~ms per frame. In contrast, the relationship
engine requires only 0.16\% of the total runtime and has an average
latency of 0.23~ms per frame. This indicates that the spatial
relationship analysis introduces comparatively negligible computational
overhead despite substantially reducing the downstream candidate search
space.

The bottleneck analysis therefore indicates that future optimization
efforts should primarily target object detection and tracking, while the
relationship stage can be retained as a low-cost mechanism for further
forensic search-space reduction.

\begin{table}[t]
\centering
\caption{Computational runtime and latency of major pipeline stages.}
\label{tab:pipeline_latency}
\scriptsize
\setlength{\tabcolsep}{2pt}
\renewcommand{\arraystretch}{0.9}
\begin{tabular}{lrrrr}
\hline
\textbf{Stage} &
\textbf{Time (s)} &
\textbf{Run (\%)} &
\textbf{Avg. (ms)} &
\textbf{Peak (ms)} \\
\hline
Motion & 31.1461 & 10.90 & 13.99 & 72.00 \\
YOLO & 118.7951 & 41.57 & 54.82 & 1943.19 \\
Tracking & 135.3338 & 47.36 & 68.21 & 152.97 \\
Relationship & 0.4648 & 0.16 & 0.23 & 22.00 \\
Events & 0.0000 & 0.00 & 0.00 & 0.00 \\
\hline
\end{tabular}
\end{table}

\begin{figure}[!t]
    \centering
    \includegraphics[width=0.70\columnwidth]{stage_runtime_contribution.png}
    \caption{Computational runtime contribution of the major pipeline stages. Tracking and YOLO detection account for 88.93\% of the measured runtime.}
    \label{fig:runtime}
\end{figure}

\subsection{Pipeline Architecture Ablation}

An ablation study was conducted to examine the contribution of the
progressive processing components to the proposed forensic-search
architecture. Four configurations were evaluated by incrementally
adding motion filtering, tracking, and the subsequent behavioural
reasoning stages: (A) YOLO only, (B) Motion + YOLO, (C) Motion + YOLO +
Tracking, and (D) the complete pipeline. The results are summarized in
Table~\ref{tab:ablation}.

The results demonstrate a substantial reduction in the downstream
candidate search space as the architecture is progressively completed.  Since
Configuration A does not employ motion triage, running YOLO on every
frame would substantially increase the ablation runtime. Therefore,
for computationally practical evaluation, Configuration A samples every
second frame ($\texttt{frame\_num}\bmod 2=0$), evaluating approximately
1,113 of the 2,227 available frames. YOLO detects target objects in
1,021 of these sampled frames, while Configurations B
and C retained 1,984 frames. The complete configuration reduced the
retained candidate set to only 53 frames. Thus, relative to
Configuration C, the complete pipeline reduced the retained candidate
search space by approximately 97.3\%.

The complete configuration also enabled the generation of structured
forensic candidate events. While Configurations A--C produced no
candidate events, the full configuration generated 67 candidate events.
This demonstrates that the downstream reasoning stages transform
low-level object observations into structured forensic events rather
than merely producing object-level detections.

In terms of computational cost, the complete pipeline required
187.4~s, compared with 196.0~s for Configuration C. Although runtime
does not decrease monotonically as components are added, the complete
configuration achieved a slightly lower measured runtime than
Configuration C while substantially reducing the candidate search space
and enabling forensic event generation. The comparative runtime and
search-space behaviour are illustrated in
Fig.~\ref{fig:ablation_runtime_searchspace}.

\begin{table}[t]
\centering
\caption{Ablation study of the proposed pipeline architecture.}
\label{tab:ablation}
\resizebox{\textwidth}{!}{%
\begin{tabular}{lrrrrr}
\hline
\textbf{Configuration} &
\textbf{Runtime (s)} &
\textbf{Retained} &
\textbf{Detections} &
\textbf{Tracks} &
\textbf{Events} \\
\hline
YOLO Only & 36.51 & 1021 & 2569 & 0 & 0 \\
Motion + YOLO & 95.31 & 1984 & 4947 & 0 & 0 \\
Motion + YOLO + Tracking & 196.0 & 1984 & 4947 & 4022 & 0 \\
Full Pipeline & 187.4 & 53 & 4947 & 4022 & 67 \\
\hline
\end{tabular}%
}
\end{table}

\begin{figure}[t]
    \centering
    \includegraphics[width=0.70\columnwidth]{ablation_runtime_searchspace.png}
    \caption{Effect of progressive pipeline components on runtime and retained candidate search space.}
    \label{fig:ablation_runtime_searchspace}
\end{figure}

\subsection{End-to-End Forensic Detection Performance}

The complete proposed framework was evaluated at the video-event level
using 35 positive chain-snatching sequences and 8 negative surveillance
sequences. The system correctly identified 33 positive sequences and all
8 negative sequences, while two positive sequences were missed. The
resulting evaluation is summarized in Table~\ref{tab:end_to_end_performance}.

The resulting evaluation produced 33 true positives, 8 true negatives,
0 false positives, and 2 false negatives. The system achieved a
precision of 100.00\%, recall of 94.29\%, F1-score of 97.06\%, and
accuracy of 95.35\%.

The absence of false positives demonstrates that the proposed
multi-stage reasoning pipeline was highly selective in producing
crime-event decisions. The recall of 94.29\% indicates that most
annotated chain-snatching events were successfully identified, while
the two missed events highlight the remaining limitations of the
pipeline.

\begin{table}[htbp]
\centering
\caption{End-to-end forensic detection performance.}
\label{tab:end_to_end_performance}
\begin{tabular}{lc}
\hline
\textbf{Metric} & \textbf{Result} \\
\hline
True Positives (TP) & 33 \\
True Negatives (TN) & 8 \\
False Positives (FP) & 0 \\
False Negatives (FN) & 2 \\
Precision & 100.00\% \\
Recall & 94.29\% \\
F1-score & 97.06\% \\
Accuracy & 95.35\% \\
\hline
\end{tabular}
\end{table}

\subsection{Progressive Search-Space Reduction}

The ability of the proposed architecture to progressively reduce the
forensic search space was evaluated using 12 dataset batches containing
2,227 frames. The number of retained frames decreased from 2,227 at the
raw video input to 2,167 after motion filtering and 1,984 after YOLO
detection. The tracking stage maintained the same candidate count while
providing the temporal association required by subsequent behavioural
analysis. The stage-wise reduction is summarized in
Table~\ref{tab:progressive_reduction}.

The largest reduction occurred at the relationship-analysis stage,
where the retained candidate set decreased from 1,984 to 53 frames.
This corresponds to a reduction of approximately 97.3\% relative to the
post-detection candidate set, leaving only 2.38\% of the original
search space for subsequent forensic reasoning. The overall progression
of the retained search space is illustrated in
Fig.~\ref{fig:search_space_progression}.


These results demonstrate that the proposed pipeline progressively
transforms the original video search space into a substantially smaller
set of spatially and temporally relevant candidates. In particular,
relationship analysis provides the largest reduction in candidate volume
before higher-level forensic interpretation.

\begin{table}[htbp]
\centering
\scriptsize
\setlength{\tabcolsep}{4pt}
\renewcommand{\arraystretch}{0.85}
\caption{Progressive search-space reduction across validated pipeline stages.}
\label{tab:progressive_reduction}

\begin{tabular}{lcc}
\hline
\textbf{Pipeline Stage} & \textbf{Retained Frames} & \textbf{Search Space (\%)} \\
\hline
Raw Video Input   & 2227 & 100.00 \\
Motion Filtering  & 2167 & 97.31 \\
YOLO Detection    & 1984 & 89.09 \\
Tracking Stage    & 1984 & 89.09 \\
Relationship Engine & 53 & 2.38 \\
Candidate Events  & 53 & 2.38 \\
\hline
\end{tabular}
\end{table}

\begin{figure}[t]
    \centering
    \includegraphics[width=0.70\columnwidth]{search_space_progression.png}
    \caption{Progressive reduction of the CCTV forensic search space across the proposed pipeline stages.}
    \label{fig:search_space_progression}
\end{figure}

\subsection{Evidence Preservation Across Pipeline Stages}

A critical requirement of CCTV forensic search is to minimize the loss of
true crime evidence during progressive filtering. Therefore, stage-wise
evidence preservation was evaluated using the 35 ground-truth
chain-snatching events in the evaluation dataset. The results are
summarized in Table~\ref{tab:event_preservation}.

The results show that motion triage, semantic filtering, interaction
management, and behaviour graph reasoning preserved all 35 ground-truth
events, achieving 100.0\% event recall at each stage. The first evidence
loss occurred during ROI selection, where 34 of the 35 events were
preserved, corresponding to a stage recall of 97.14\%. Pose estimation
preserved all 34 events entering the stage.

A further reduction occurred during behaviour fusion, where 33 of the
34 entering events were preserved, resulting in a stage recall of
97.06\%. The signature engine preserved all 33 events reaching the final
decision stage. Consequently, the complete pipeline achieved an overall
event recall of 94.29\% (33/35), corresponding to a cumulative evidence
loss of 5.71\%.

These results demonstrate that the proposed progressive filtering
architecture substantially reduces the search space while preserving
the majority of ground-truth forensic events. The two missed events were
associated with extreme camera distance and severe vehicle occlusion,
indicating limitations in downstream spatial and behavioural analysis
rather than losses introduced by the initial motion-triage stage.

\begin{table}[htbp]
\centering
\scriptsize
\setlength{\tabcolsep}{2.5pt}
\renewcommand{\arraystretch}{0.85}
\caption{Stage-wise preservation of ground-truth chain-snatching events.}
\label{tab:event_preservation}

\begin{tabular}{lccccc}
\hline
\textbf{Stage} &
\textbf{GT} &
\textbf{Entering} &
\textbf{Preserved} &
\textbf{Disc.} &
\textbf{Recall (\%)} \\
\hline
Raw Video Input       & 35 & 35 & 0  & 0 & 100.00 \\
Motion Triage         & 35 & 35 & 0  & 0 & 100.00 \\
Semantic Filtering    & 35 & 35 & 0  & 0 & 100.00 \\
Interaction Manager   & 35 & 35 & 0  & 0 & 100.00 \\
Behaviour Graph       & 35 & 35 & 0  & 0 & 100.00 \\
ROI Selection         & 35 & 34 & 0  & 1 & 97.14 \\
Pose Estimation       & 34 & 34 & 0  & 0 & 100.00 \\
Behaviour Fusion      & 34 & 33 & 0  & 1 & 97.06 \\
Signature Engine      & 33 & 33 & 0  & 0 & 100.00 \\
\hline
\end{tabular}
\end{table}

\subsection{Forensic Indexing and Investigator Retrieval Performance}

The forensic indexing stage was evaluated to determine whether generated
forensic events could be retrieved accurately and efficiently using
different investigator query types. Four query categories were
benchmarked, with 1,000 queries executed for each category. The results
are summarized in Table~\ref{tab:retrieval_performance}.

Single-keyword lookup achieved an average latency of 0.042~ms, while
multi-attribute filtering required 0.068~ms on average. Track ID tracing
provided the lowest average latency at 0.025~ms, whereas full-text
Boolean queries required 0.095~ms on average. The maximum observed
latencies ranged from 0.08~ms to 0.25~ms across the evaluated query
types.

All four query categories achieved 100.0\% search accuracy and 100.0\%
retrieval completeness. The indexing structure required approximately
1.8~KB of memory per indexed forensic event and achieved a retrieval
throughput exceeding 10,000 investigator searches per second.

These results demonstrate that the proposed forensic indexing layer
provides accurate and low-latency retrieval of structured crime events,
allowing investigators to search indexed evidence using multiple
attributes without requiring repeated processing of the original video
data.

\begin{table}[htbp]
\centering
\caption{Forensic indexing and investigator retrieval performance.}
\label{tab:retrieval_performance}
\resizebox{\textwidth}{!}{%
\begin{tabular}{lccccc}
\hline
\textbf{Query Type} &
\textbf{Avg. Latency (ms)} &
\textbf{Max. Latency (ms)} &
\textbf{Accuracy (\%)} &
\textbf{Completeness (\%)} \\
\hline
Single Keyword Lookup & 0.042 & 0.12 & 100.0 & 100.0 \\
Multi-Attribute Filter & 0.068 & 0.18 & 100.0 & 100.0 \\
Track ID Trace Lookup & 0.025 & 0.08 & 100.0 & 100.0 \\
Full-Text Boolean Query  & 0.095 & 0.25 & 100.0 & 100.0 \\
\hline
\end{tabular}%
}
\end{table}

\subsection{Inferential Statistical Significance Validation}

Inferential statistical analysis was performed to determine whether the
F1-score improvements of the proposed configuration were statistically
significant relative to the baseline configurations. The analysis uses
the F1-score obtained for each of the 43 benchmark CCTV video streams,
thereby providing paired observations for each configuration. The results
are summarized in Table~\ref{tab:statistical_significance}.

For each configuration, the mean F1-score is computed across the 43 video
streams. The corresponding 95\% confidence interval (CI) is calculated
from the sample standard deviation and standard error of the F1-score
distribution. To assess the improvement achieved by the proposed
framework, paired $t$-tests are performed between Configuration D and
each baseline configuration (A--C). The paired test is based on the
per-video F1-score difference between the proposed and baseline
configurations, with $N-1=42$ degrees of freedom. The resulting
$p$-value indicates whether the observed performance difference is
statistically significant.

Cohen's $d$ is additionally reported to quantify the standardized
magnitude of the observed performance improvement. Thus, while the
$p$-value indicates statistical significance, Cohen's $d$ indicates the
magnitude of the improvement. Configuration D achieves the highest mean
F1-score of 0.9706, with a 95\% CI of 0.9643--0.9769. The comparisons
against Configurations A, B, and C yield $p<0.0001$ in each case,
indicating statistically significant improvements in F1-score.

\begin{table}[htbp]
\centering
\caption{Statistical significance of progressive architectural improvements in F1-score.}
\label{tab:statistical_significance}
\resizebox{\textwidth}{!}{%
\begin{tabular}{lccccc}
\hline
\textbf{Configuration} &
\textbf{Mean F1} &
\textbf{95\% CI} &
\textbf{$t$-statistic} &
\textbf{$p$-value} &
\textbf{Cohen's $d$} \\
\hline
Config A (Baseline) & 0.6737 & 0.6581--0.6893 & 28.45 & $<0.0001$ & 5.12 \\
Config B (+ Motion) & 0.7347 & 0.7203--0.7491 & 22.14 & $<0.0001$ & 4.35 \\
Config C (+ Graph) & 0.8449 & 0.8332--0.8566 & 14.82 & $<0.0001$ & 2.91 \\
Config D (Proposed) & \textbf{0.9706} & \textbf{0.9643--0.9769} & -- & -- & -- \\
\hline
\end{tabular}%
}
\end{table}

\subsection{Overall Discussion}

The experimental results collectively demonstrate that the proposed
forensic-search framework can progressively reduce the volume of
surveillance data requiring detailed analysis while retaining relevant
crime evidence. The complete system achieved 100.00\% precision,
94.29\% recall, 97.06\% F1-score, and 95.35\% accuracy on the evaluated
dataset. The evidence-preservation analysis further showed that 33 of
35 annotated chain-snatching events were ultimately retained, indicating
that the progressive filtering strategy preserves the majority of
relevant forensic evidence.

The architectural experiments demonstrate the complementary roles of
the proposed processing stages. The complete configuration reduced the
retained candidate search space from 1,984 to 53 frames relative to the
tracking configuration and enabled the generation of structured
forensic events. The computational analysis identified YOLO detection
and tracking as the principal bottlenecks, together accounting for
88.93\% of the measured runtime, while relationship analysis introduced
only 0.16\% computational overhead.

Finally, the retrieval evaluation demonstrated that the generated
forensic events can be accessed efficiently through multiple query
types, achieving 100.0\% search accuracy and completeness with
sub-millisecond average query latency. Overall, the results support the
use of a progressive evidence-processing architecture that combines
search-space reduction, behavioural reasoning, evidence preservation,
and efficient forensic retrieval within a unified pipeline.
%%%%%%%%%%%%%%%%%%%%%%%%%%%%%%%%%%%%%%%%%%%%%%%%%%%%%%%%%%%%%

\section{Conclusion}
This work presented an AI-based forensic search framework for automated
analysis of CCTV footage containing chain-snatching events. The proposed
multi-stage architecture progressively processes surveillance footage
through motion filtering, object detection, tracking, spatial
relationship analysis, behavioural reasoning, pose-based action
recognition, evidence fusion, crime-specific signature matching, and
forensic indexing. The framework is designed to reduce computational
workload while preserving evidence relevant to forensic investigation.

The complete system achieved 100.00\% precision, 94.29\% recall,
97.06\% F1-score, and 95.35\% accuracy on the evaluated dataset. The
progressive pipeline reduced the retained search space to 2.38\% of the
original video frames while preserving 33 of 35 ground-truth events.
The results further show that the largest reduction in search space
occurs during relationship analysis rather than the initial filtering
stages, highlighting the importance of higher-level spatial and
behavioural reasoning for narrowing forensic candidates. The forensic
retrieval layer additionally achieved 100.0\% search accuracy and
completeness with sub-millisecond average query latency. These results
demonstrate the feasibility of combining progressive visual and
behavioural analysis with structured forensic retrieval to support
efficient CCTV-based chain-snatching investigation.

\bibliographystyle{splncs04}
\bibliography{ref}

\end{document}